%%%%%%%%%%%%%%%%%%%%%%%%%%%%%%%%%%%%%%%%%%%%%%%%%%%%%%%%%%%%%%%%%%%%%%%%%%%%%%%%
%% 
\cleardoubleoddpage%  Make sure to start each chapter on a new odd page
\chapter{Evaluation Criteria}
\label{sec:evaluation}
This chapter defines the criteria used to evaluate the selected static analysis tools. The criteria cover both practical aspects, such as installation and usability, 
and technical aspects, such as format support and the quality of the generated analysis output. Since the selected tools target different executable formats, 
comparisons are performed both across the complete tool set and within the respective categories of native-code and managed-code analysis tools.

\section{Installation and Setup}
\label{sec:evaluation:installation}
The installation process is evaluated from the perspective of a user who obtains the tools directly from their official sources. For this purpose, only precompiled binaries, 
installation packages, or installers provided by the respective developers are used. The tools are neither compiled from source nor installed through the package repositories 
of Debian or another Linux distribution.

This approach avoids differences introduced by third-party package maintainers, such as outdated versions, distribution-specific modifications, or additional packaging decisions. 
It also ensures that the evaluated installation procedure corresponds as closely as possible to the method recommended or directly supported by each project. 
Compiling the tools from source is excluded because the associated build requirements and potential compilation problems would evaluate the development process rather than the 
installation experience intended for regular users.

The evaluation considers the availability of suitable Linux packages, the number and complexity of the required installation steps, additional dependencies, necessary configuration, 
and any problems encountered before the tool can be used for its first analysis. The source and version of each installed tool are recorded to make the evaluation reproducible.

The latest stable version available at the time of the evaluation is selected for each tool, while nightly builds and other development versions are excluded. 
Unless required for basic operation, the tools are evaluated using their default configuration and without optional plugins or extensions. Account creation, 
licence activation, and first-start configuration are considered part of the setup process. Installation is regarded as complete once the application can be started successfully 
and an appropriate test file can be imported for analysis.

\section{User Interface and Analysis Workflow}
\label{sec:evaluation:uiux}
The user interface and analysis workflow are evaluated by performing a set of common static analysis tasks in each applicable tool. These tasks include importing a file, 
initiating the analysis, locating functions and strings, following cross-references, and navigating between related parts of the program. Where supported, 
the accessibility of disassembly, decompilation, and graphical control-flow representations is also considered.

The evaluation further examines how clearly the tools present the available information and how easily the analyst can move between different views. Basic interaction features, 
such as renaming functions and variables or adding comments, are considered where they form part of the intended workflow. Particular attention is given to whether important functionality 
is readily discoverable or requires prior knowledge of the tool.

Since usability cannot be measured entirely objectively, all tools are evaluated by the same "analyst" using comparable tasks. The assessment is therefore based on the consistency and 
practicality of the workflow rather than on personal preferences concerning visual appearance.

\section{Supported File Formats and Architectures}
\label{sec:evaluation:fileFormats}
The supported file formats and processor architectures are determined from the official documentation of each tool (if there is any). For the formats represented in the test dataset, 
support is additionally verified by importing and analyzing the corresponding files. This distinction separates officially stated capabilities from those confirmed during the 
practical evaluation. The comparison focuses on formats relevant to this work, particularly native executables, .NET assemblies, and Java bytecode.

\section{Quality of Analysis Output}
\label{sec:evaluation:output}
The quality of the generated analysis output is evaluated with regard to both correctness and readability. Relevant aspects include the identification of functions and function calls, 
the reconstruction of control flow, the representation of variables and data types, and the preservation of the original program logic. Incorrectly identified code, missing information, 
and output that could mislead the analyst are also considered.

For the synthetic samples, the generated disassembly and decompiled code are compared with the corresponding source code. This provides a known reference against which the reconstructed 
program logic can be examined. Since the original source code of the malware samples is not available, their output is instead assessed according to whether it supports the 
identification and understanding of relevant program behaviour.

Because native machine code, .NET assemblies, and Java bytecode contain different amounts and forms of recoverable information, output quality is primarily compared between tools that 
analyze the same type of program. Screenshots and selected code excerpts are used in the practical evaluation to illustrate significant differences between the tools.


%% 
%%%%%%%%%%%%%%%%%%%%%%%%%%%%%%%%%%%%%%%%%%%%%%%%%%%%%%%%%%%%%%%%%%%%%%%%%%%%%%%%
